\BourbakiExercisesSection

\begin{enumerate}
\def\labelenumi{\arabic{enumi}.}
\item
  Let \(\top\) be a law of composition on a set E. Denoting by F the sum set (Set Theory, II, § 4, no. 8) of E and a set \(\{e\}\) consisting of a single element and identifying E and \(\{e\}\) with the corresponding subsets of F, show that it is possible in one and only one way to define on F a law of composition \(\overline{\top}\) which induces on E the law \(\top\) and under which \(e\) is identity element; if \(\top\) is associative, the law \(\overline{\top}\) is associative. (F is said to be derived from E ``by adjoining an identity element''.)
\item
  Let \(\mathsf{T}\) be a law everywhere defined on E.
\end{enumerate}

\begin{enumerate}
\def\labelenumi{(\alph{enumi})}
\item
  For \(\top\) to be associative, it is necessary and sufficient that every left translation \(\gamma_{x}\) be permutable with every right translation \(\delta_{y}\) in the set of mappings of E into E (with the law \(f \circ g\) ).
\item
  Suppose that E has an identity element. For T to be associative and commutative, it is necessary and sufficient that the mapping \((x,y)\mapsto x\top y\) be a morphism of the magma \(E\times E\) into the magma E.
\end{enumerate}

\begin{enumerate}
\def\labelenumi{\arabic{enumi}.}
\setcounter{enumi}{2}
\item
  In the set F of mappings of E into E, for the relation \(f \circ g = f \circ h\) to imply g = h, it is necessary and sufficient that f be injective; for the relation \(g \circ f = h \circ f\) to imply g = h, it is necessary and sufficient that f be surjective; for f to be cancellable (under the law \(\circ\) ), it is necessary and sufficient that f be bijective.
\item
  For there to exist on a set E a law of composition such that every permutation of E is an isomorphism of E onto itself under this law, it is necessary and sufficient that E have 0, 1 or 3 elements.
\item
  For \(2 \leqslant n \leqslant 5\) , determine on a set E with n elements all the laws everywhere defined admitting an identity element and under which all the elements of E are cancellable; for n = 5 show that there exist non-associative laws satisfying these conditions.
\end{enumerate}

(N.B. Exercises 6 to 17 inclusive refer to associative laws written multiplicatively; \(e\) denotes the identity when it exists, \(\gamma_a\) and \(\delta_a\) the translations by \(a \in E\) ; for every subset \(X\) of \(E\) we write \(\gamma_a(X) = aX\) , \(\delta_a(X) = Xa\) .)

\begin{enumerate}
\def\labelenumi{\arabic{enumi}.}
\setcounter{enumi}{5}
\item
  For an associative law on a finite set, every cancellable element is invertible (use no. 3, Proposition 3).
\item
  Given an associative law on a set E and an element \(x \in E\) , let A be the set of \(x^{n}\) for \(n \in N^{*}\) ; if there is an identity element, let B be the set of \(x^{n}\) for \(n \in N\) ; if further x is invertible, let C be the set of \(x^{a}\) for \(a \in Z\) . Show that, if A (resp.
\end{enumerate}

B, C) is infinite, it is isomorphic (with the law induced by the given law on E) to \(N^{*}\) (resp. N, Z) with addition.

\begin{enumerate}
\def\labelenumi{\arabic{enumi}.}
\setcounter{enumi}{7}
\item
  In the notation of Exercise 7, suppose A is finite; show that A contains one and only one idempotent (§ 1, no. 4) (observe that if \(x^p\) and \(x^q\) are idempotents, then \(x^p = x^{pq}\) , \(x^q = x^{pq}\) , hence \(x^p = x^q\) ); if \(h = x^p\) , the set of \(x^n\) for \(n \geqslant p\) is a stable subset D of E such that under the law induced on D, D is a group.
\item
  Under a multiplicative law on a set E let \(a\) be a left cancellable element of E.
\end{enumerate}

\begin{enumerate}
\def\labelenumi{(\alph{enumi})}
\item
  If there exists an element \(u\) such that \(au = a\) , show that \(ux = x\) for all \(x \in \mathbb{E}\) ; in particular, if \(xu = x\) for all \(x \in \mathbb{E}\) , \(u\) is identity element.
\item
  If there exist \(u \in E\) such that au = a and \(b \in E\) such that ab = u, show that ba = u (form aba); in particular, if there exists an identity element e and an element b such that ab = e, b is the inverse of a.
\end{enumerate}

\begin{enumerate}
\def\labelenumi{\arabic{enumi}.}
\setcounter{enumi}{9}
\tightlist
\item
  If a and b are two elements of E such that ba is left cancellable, show that a is left cancellable. Deduce that under an associative and commutative law on E the set S of non-cancellable elements of E is such that ES ⊆ S (and in particular is stable).
\end{enumerate}

¶ 11. E is called a left semi-group if every element x of E is left cancellable. (a) If u is an idempotent (§ 1, no. 4), then ux = x for all \(x \in E\) (use Exercise 9(a)); u is identity element for the law induced on Eu.

\begin{enumerate}
\def\labelenumi{(\alph{enumi})}
\setcounter{enumi}{1}
\item
  If \(u\) and \(v\) are two distinct idempotents of \(\mathbf{E}\) , then \(\mathrm{Eu} \cap \mathrm{Ev} = \varnothing\) and the sets \(\mathrm{Eu}\) and \(\mathrm{Ev}\) (with the laws induced by that on \(\mathbf{E}\) ) are isomorphic.
\item
  Let R be the complement of the union of the sets Eu, where u runs through the set of idempotents of E. Show that ER ⊂ R; it follows that R is a stable subset of E. If R is non-empty, then \(aR \neq R\) for all \(a \in R\) (use Exercise 9(a) to prove that, in the contrary case, R would contain an idempotent); in particular, R is then infinite. R is called the residue of the left semi-group E.
\item
  If R is non-empty and there exists at least one idempotent u in E (that is \(E \neq R\) ), for all \(x \in REu\) , \(xEu \neq Eu\) ; in particular, no element of REu is invertible in Eu and REu is an infinite set (use Exercise 8).
\item
  If E has an identity element e, e is the only idempotent in E and R is empty (note that E = Ee).
\item
  If there exists a right cancellable element \(a\) in E (in particular if the given law on E is commutative), either E has an identity element or \(\mathbf{E} = \mathbf{R}\) (note that if there exists an idempotent \(u\) , then \(xa = xua\) for all \(x \in \mathbf{E}\) ). For example, if E is the set of integers \(\geqslant 1\) , with addition, E is a commutative semi-group such that \(\mathbf{E} = \mathbf{R}\) .
\item
  If there exists \(a \in \mathbf{E}\) such that \(\delta_a\) is surjective, either \(\mathbf{E}\) has an identity element or \(\mathbf{E} = \mathbf{R}\) (examine separately the case where \(a \in \mathbf{R}\) and the case where \(a \in \mathbf{Eu}\) for an idempotent \(u\) ).
\end{enumerate}

¶ 12. Under a multiplicative law on E, let \(a \in E\) be such that \(\gamma_{a}\) is surjective.

\begin{enumerate}
\def\labelenumi{(\alph{enumi})}
\item
  Show that, if there exists \(u\) such that \(ua = a\) , then \(ux = x\) for all \(x \in E\) .
\item
  For an element \(b \in E\) to be such that ba is left cancellable, it is necessary and sufficient that \(\gamma_{a}\) be surjective and that b be left cancellable.
\end{enumerate}

¶ 13. Suppose that, for all \(x \in E\) , \(\gamma_x\) is surjective. Show that if, for an element \(a \in E\) , \(\gamma_a\) is a bijective mapping, \(\gamma_x\) is a bijective mapping for all \(x \in E\) (use Exercise 12(b)); this holds in particular if there exist two elements \(a, b\) of \(E\) such that \(ab = b\) (use Exercise 12(a)); \(E\) is then a left semi-group whose residue \(R\) is empty; moreover, for every idempotent \(u\) , every element of \(Eu\) is invertible in \(Eu\) .

¶ 14. For an associative law on a finite set E there exist minimal subsets of the form aE (that is minimal elements of the set of subsets of E of this form, ordered by inclusion).

\begin{enumerate}
\def\labelenumi{(\alph{enumi})}
\item
  If M = aE is minimal, then xM = xE = M for all \(x \in M\) ; with the law induced by that on E, M is a left semi-group (Exercise 11) in which every left translation is bijective (cf.~Exercise 13).
\item
  If \(\mathbf{M} = a\mathbf{E}\) and \(\mathbf{M}' = a'\mathbf{E}\) are minimal and distinct, then \(\mathbf{M} \cap \mathbf{M}' = \varnothing\) ; for all \(b \in \mathbf{M}\) , the mapping \(x' \mapsto bx'\) of \(\mathbf{M}'\) into \(\mathbf{M}\) is a bijective mapping of \(\mathbf{M}'\) onto \(\mathbf{M}\) . Deduce that there exist an idempotent \(u' \in \mathbf{M}'\) such that \(bu' = b\) and an idempotent \(u \in \mathbf{M}\) such that \(uu' = u\) (take \(u\) to be the idempotent such that \(bu = b\) ). Show that \(u'u = u'\) (consider \(uu'u\) ) and that every \(y' \in \mathbf{M}'\) such that \(y'u = y'\) belongs to \(\mathbf{M}'u'\) .
\item
  Show that the mapping \(x' \mapsto ux'\) of \(M'u'\) into \(M\) is an isomorphism of \(M'u'\) onto \(Mu\) ; deduce that \(M\) and \(M'\) are isomorphic left semi-groups.
\item
  Let \(M_{i}\) ( \(1 \leqslant i \leqslant r\) ) be the distinct minimal subsets of the form aE: deduce from (b) that the idempotents of each left semi-group \(M_{i}\) can be arranged in a sequence \((u_{ij})\) ( \(1 \leqslant j \leqslant s\) ) so that \(u_{ij}u_{kj} = u_{ij}\) for all i, j, k. If K is the union of the \(M_{i}\) , show that \(Eu_{ij} \subset K\) (note that, for all \(x \in E\) , \(xu_{ij}E\) is a minimal subset of the form aE); deduce that \(Eu_{ij}\) is the union of the \(u_{kj}Eu_{kj}\) for \(1 \leqslant k \leqslant r\) (show that \((\mathrm{Eu}_{ij}) \cap \mathbf{M}_{k} = u_{kj}\mathrm{Eu}_{kj}\) ) and that \(Eu_{ij}E = K\) . Prove finally that every minimal subset of the form Eb is identical with one of the s sets \(Eu_{ij}\) (note, using (a)), that \(Ebu_{ij}b = Eb\) and deduce that \(Eb \subset K\) ).
\end{enumerate}

\begin{enumerate}
\def\labelenumi{\arabic{enumi}.}
\setcounter{enumi}{14}
\tightlist
\item
  Let \((x_{i})_{1\leqslant i\leqslant n}\) be a finite sequence of left cancellable elements.
\end{enumerate}

\begin{enumerate}
\def\labelenumi{(\alph{enumi})}
\item
  Show that the relation \(x_{1}x_{2}\ldots x_{n} = e\) implies all the relations \(x_{i + 1}\ldots x_nx_1x_2\ldots x_i = e\) which are obtained from one another by ``cyclic permutation'' for \(1\leqslant i\leqslant n\) .
\item
  Deduce that if the composition of the sequence \((x_{i})\) is invertible each of the \(x_{i}\) is invertible.
\end{enumerate}

¶ 16. Let \(\top\) be an associative law on a set E and E* the set of cancellable elements of E; suppose that E* is non-empty and that every cancellable element is a central element. Let \(\mathfrak{F}\) denote the set of subsets \(\mathbf{X}\) of \(\mathbf{E}\) with the following property: there exists \(y \in \mathrm{E}^*\) such that \(\delta_y(\mathrm{E}) \subset \mathbf{X}\) .

\begin{enumerate}
\def\labelenumi{(\alph{enumi})}
\item
  Show that the intersection of two sets of \(\mathfrak{F}\) belong to \(\mathfrak{F}\) .
\item
  Let \(\Phi\) be the set of functions \(f\) defined on a set of \(\mathfrak{F}\) taking their values in
\end{enumerate}

E and such that, for \(X \in \mathfrak{F}\) , \(f(\mathbf{X})\) belongs to F. Let R denote the relation between elements f and g of \(\Phi\) : ``there exists a set \(X \in F\) such that the restrictions of f and g to X are identical''. Show that R is an equivalence relation; let \(\Psi = \Phi/R\) be the quotient set of \(\Phi\) under this relation.

\begin{enumerate}
\def\labelenumi{(\alph{enumi})}
\setcounter{enumi}{2}
\item
  Let \(f\) and \(g\) be two elements of \(\Phi\) , \(A \in \mathfrak{F}\) and \(B \in \mathfrak{F}\) the sets where \(f\) and \(g\) are respectively defined; there exists \(X \subset B\) and belonging to \(\mathfrak{F}\) such that \(g(X) \subset A\) ; if \(g_X\) is the restriction of \(g\) to \(X\) , show that the mapping \(f \circ g_X\) belongs to \(\Phi\) and that its class (mod. R) depends only on the classes of \(f\) and \(g\) (and not on \(X\) ); this class is called the composition of that of \(f\) and that of \(g\) ; show that the law of composition thus defined on \(\Psi\) is associative and has an identity element.
\item
  For all \(a \in \mathbb{E}\) , show that the left translation \(\gamma_{a}\) belongs to \(\Phi\) ; let \(\phi_{a}\) be its class (mod. R). Show that the mapping \(x \mapsto \phi_{x}\) is an isomorphism of E onto a submonoid of \(\Psi\) and that, if \(x \in \mathbb{E}^{*}\) , \(\phi_{x}\) is invertible in \(\Psi\) (consider the inverse mapping of \(\gamma_{x}\) , show that it belongs to \(\Phi\) and that its class (mod. R) is the inverse of \(\phi_{x}\) ). Deduce a generalization of Theorem 1 of no. 4.
\end{enumerate}

\begin{enumerate}
\def\labelenumi{\arabic{enumi}.}
\setcounter{enumi}{16}
\tightlist
\item
  Let E be a monoid and S a stable subset of E with the following properties:
\end{enumerate}

\((\alpha)\) For all \(s \in S\) and all \(a \in E\) , there exist \(t \in S\) and \(b \in E\) such that \(sb = at\) .

(β) For all \(a, b \in \mathbb{E}\) and all \(s \in \mathbb{S}\) such that \(sa = sb\) , there exist \(t \in \mathbb{S}\) such that \(at = bt\) .

\begin{enumerate}
\def\labelenumi{(\alph{enumi})}
\tightlist
\item
  In E × S let \((a, s) \sim (b, t)\) denote the relation: ``there exist \(s'\) and \(t'\) in S such that \(tt' = ss'\) and \(as' = bt'\) .''
\end{enumerate}

Show that \(\sim\) is an equivalence relation. We write \(\overline{\mathbf{E}} = \mathbf{E} \times \mathbf{S}/\sim\) , denote by \(as^{-1}\) the equivalence class of \((a, s) (\text{mod } \sim)\) and by \(\varepsilon: \mathbf{E} \to \overline{\mathbf{E}}\) the mapping \(a \mapsto ae^{-1}\) .

\begin{enumerate}
\def\labelenumi{(\alph{enumi})}
\setcounter{enumi}{1}
\item
  Show that there exists on \(\overline{\mathbf{E}}\) one and only one monoid structure such that \(\varepsilon\) is a unital homomorphism and such that, for all \(s \in \mathbf{S}\) , \(\varepsilon(s)\) is invertible.
\item
  Show that \((\overline{\mathbf{E}},\varepsilon)\) has the universal property described in Theorem 1 (no. 4).
\item
  Show that \(\varepsilon\) is injective if and only if the elements of S are cancellable.
\end{enumerate}

